\chapter{复数}

\section{复数及其四则运算}
\subsection{复数的概念和四则运算}
\zsd 复数的概念和四则运算 \\
A组：
\begin{enumerate}
    
\item 若复数$z=\dfrac{m^2+m-2}{m-3}+(m^2-4m+3)\mi$是纯虚数, 则实数$m=$\blkda{$-2$}. 

\item 设复数集$\bC$为全集, 若集合$A=\{z\mid z=a\mi, a\zr\}$, 集合$B=\{z\mid z=a+b\mi, a\zr, a\neq 0, b\zr\}$, 
则$\complement_{\bC}(A\cup B)=$\blkda{$\varnothing$}. 

\item 下列命题中, 不正确命题的序号是\blkda{\ding{173}\ding{174}\ding{175}}: 
\begin{dingautolist}{172}
\item 对任意两个复数$x,y$, 若满足$x>y$, 则$x,y$必定都是实数; 
\item 复数$z=a+b\mi(a,b\zr)$的虚部是$b\mi$; 
\item 当$a=0$时, 复数$z=a+b\mi(a,b\zr)$为纯虚数; 
\item 因为$\mi$表示虚数单位, 所以它不是一个虚数. 
\end{dingautolist}

\item $a=0$ 是复数 $a+b i(a, b \in R)$ 为纯虚数的（\blkda[3]{必要不充分}）条件．

\item 复数 $z=-2-i$ 的虚部是\blkda{-1} ．

\item 在复数范围内，-4 的所有平方根为\blkda{$\pm 2i$} ．

\end{enumerate}


B组：
\begin{enumerate}
    \item 设$x,y\zr$, 且$\dfrac{x}{1-\mi}+\dfrac{y}{1-2\mi}=\dfrac{5}{1-3\mi}$, 则$x+y=$\blkda{4}. 

    \item 已知$z_1=m$, $z_2=2-\mi$, 且$\abs{z_1}>\abs{z_2}$, 则实数$m$的取值范围是\blkda[4]{$(-\infty,-\sqrt{5})\cup(\sqrt{5},+\infty)$}. 

    \item 已知$z_1=(a^2-a-6)+(1-2a)\mi$, $z_2=(a-3)+(a^2-2a+2)\mi$, 其中$a\zr$. 若$\overline{z_1}=z_2$, 则$a=$\blkda{3}. 

    \item 以下四个式子中, 所有正确的式子的序号是\blkda{\ding{174}\ding{175}}: 
\begin{dingautolist}{172}
\item $\abs{a+b\mi}>0(a,b\zr)$; 
\item $4z>3z$, 其中$\abs{z}>2$; 
\item $\abs{a+b\mi}=\abs{b-a\mi}(a,b\zr)$; 
\item 若$a\zc$, 则$\sqrt{\abs{a^2}}=\abs{a}$. 
\end{dingautolist}

\item 复数$z=a^2-b^2+(a+\abs{b})\mi(a,b\zr)$为纯虚数的充要条件是\dotfill(\kh{D})
\xx{$a=\pm b$}{$a<0$且$a=-b$}{$a>0$且$a=b$}{$a>0$且$a=\pm b$}

\item 以下各命题中, 正确的命题为\dotfill(\kh{D})
\xx
{已知两个复数$z_1,z_2$, 若$z_1-z_2>0$, 则$z_1>z_2$}
{两个共轭复数的差是纯虚数}
{若$z$为复数, 且$m,n$都是正有理数, 则$(z^m)^n=z^{mn}$}
{若$z\zc$, 且$(z+\mi)^2=\abs{z+\mi}^2$, 则$z\notin\bR$}


\item ``两个复数的和是虚数''的一个必要非充分条件是\dotfill(\kh{C})
\xx
{两个数都是虚数}
{两个数中一个是虚数, 一个是实数}
{ 两个数中至少有一个是虚数}
{两个数都是纯虚数}

\item $z_1=\overline{z_2}$一个必要非充分条件是\dotfill(\kh{B})
\xx{$\abs{z_1-\overline{z_2}}=0$}
{$\abs{z_1}=\abs{z_2}$}
{$z_1=z_2$}
{$\overline{z_1}=z_2$}

\item 对任意复数$z=x+y\mi(x,y\zr)$, 下列结论正确的是\dotfill(\kh{D})
\xx{$\abs{z-\overline{z}}=2y$}
{$z^2=x^2+y^2$}
{$\abs{z-\overline{z}}\geq 2x$}
{$\abs{z}\leq\abs{x}+\abs{y}$}

\item 设复数 $z=a-i$ ，其中 $i$ 为虚数单位，$a \in R$ ．（1）若 $z(1+i)$ 是纯虚数，求实数 $a$ 的值；（2）若 $a=2$ ，求复数 $\frac{z}{1+i}+i$ 的模．
\begin{jieda}
    (1) $a=-1$;(2)$\sqrt{2}/2$
\end{jieda}

\item 已知复数 $z$ 是纯虚数，$(z+2)^2-8 i$ 是实数．
（1）求 $z$ ；
（2）若 $\frac{1}{z_1}=\frac{1}{z+2}-z$ ，求 $\left|z_1\right|$ ．
\begin{jieda}
    (1)$2i$; (2)$\dfrac{2\sqrt{82}}{41}$
\end{jieda}

\item 已知复数$z_1,z_2$满足$\abs{z_1}=\abs{z_2}=1$, 且$z_1+z_2=\mi$, 分别求$z_1,z_2$. 
\begin{jieda}
    当$z_1=\dfrac{\sqrt{3}}{2}+\dfrac{1}{2}i$时，$z_2=-\dfrac{\sqrt{3}}{2}+\dfrac{1}{2}i$;
     当$z_1=-\dfrac{\sqrt{3}}{2}+\dfrac{1}{2}i$时，$z_2=\dfrac{\sqrt{3}}{2}+\dfrac{1}{2}i$
\end{jieda}

\item 在复平面上作出满足下列条件的复数在复平面上对应的点集所表示的图形: 
\begin{enumerate}[(1)]
\item $1\leq\abs{z}<3$; 
\item $1<\Re{z}<2$, 且$1<\Im{z}<2$; 
\item $\abs{z}>3$, 且$\Im{z}<-1$; 
\item $\abs{z}=2$, 且$\frac{1}{2}<\Im{z}<2$. 
\end{enumerate}

\item 已知复数$z=x+y\mi(x,y\zr)$满足$\abs{z-4\mi}=\abs{z+2}$, 求$2^x+4^y$的最小值以及相应的$x,y$的值. 
\begin{jieda}
    当$x=\dfrac{3}{2}$时，有最小值$4\sqrt{2}$
\end{jieda}

\item 已知复数 $z_1=m+\left(4-m^2\right) i$ ，复数 $z_2=2 \cos \theta+(\lambda+3 \sin \theta) i$ ，其中 $m 、 \lambda \in \mathbf{R}$ ，且 $z_1=z_2$ ．求 $\lambda$ 取值范围．
\begin{jieda}
    $\lambda \in [-\dfrac{9}{16},7]$
\end{jieda}
\end{enumerate}

\section{复数及其四则运算}
\subsection{复数的概念和四则运算}
\zsd 复数的概念和四则运算 \\

\begin{enumerate}
    \item 已知复数$z=\dfrac{1+3\mi}{1-\mi}$, $\mi$为虚数单位, 则$\abs{z}$为\blkda{$\sqrt{5}$}. 


\item 计算: $\mi\cdot \mi^3\cdot \mi^5\cdot \cdots \cdot \mi^{2005}=$\blkda{$\mi$}. 


\item 计算: $(1+\mi)^{10}-(1-\mi)^{10}=$\blkda{$64\mi$}. 


\item 已知$f(z)=\overline{z}-1$, 且$f(z_1-z_2)=4+4\mi$. 若$z_1=2-3\mi$, 则复数$z_2=$\blkda{$-3+\mi$}, $\abs{z_1+z_2}=$\blkda{$\sqrt{5}$}. 


\item 已知$z\zc$, 且$\dfrac{z-1}{z+1}$是纯虚数, 则复数$z$对应的点的轨迹在复平面上所对应的直角坐标方程是\blkda[4]{$x^2+y^2=1(x\neq\pm1)$}. 


\item 设$z$是复数, $a(z)$表示满足$(1+z)^n$是正实数的最小正整数$n$, 则对虚数单位$\mi$, $a(\mi)=$\dotfill(\kh{A})
\xx{8}{6}{4}{2}


\item 已知复数$z_1=\cos\alpha+\mi\sin\alpha$, $z_2=\cos\beta+\mi\sin\beta$, 则复数$z_{1}\cdot z_{2}$的实部是\dotfill(\kh{D})
\xx
{$\sin(\alpha-\beta)$}
{$\sin(\alpha+\beta)$}
{$\cos(\alpha-\beta)$}
{$\cos(\alpha+\beta)$}


\item 若$\left(\dfrac{1+\mi}{1-\mi}\right)^n=\mi$, 则整数$n$等于\dotfill(\kh{B})
\xx
{$4k$(其中$k\zz)$}
{$4k+1$(其中$k\zz)$}
{$4k+2$(其中$k\zz)$}
{$4k+3$(其中$k\zz)$}


\item 已知$x,y\zr$, 若$P=2(x+y\mi)(x-y\mi)$, $Q=\abs{2\sqrt{xy}+(x-y)\mi}^2$, 则$P,Q$的大小关系是\dotfill(\kh{B})
\xx{$P>Q$}{$P\geq Q$}{$P\leq Q$}{不能确定}


\item $\triangle{ABC}$的三个顶点所对应的复数分别为$z_1,z_2,z_3$. 若复数$z$满足$\abs{z-z_1}=\abs{z-z_2}=\abs{z-z_3}$, 则$z$所对应的点是$\triangle{ABC}$的\dotfill(\kh{B})
\xx{内心}{外心}{重心}{垂心}

\item 已知$\omega=-\dfrac{1}{2}+\dfrac{\sqrt{3}}{2}\mi$, 求证: 
\begin{enumerate*}[(1)]
\item $\omega^3=1$; 
\item $\omega^2=\overline{\omega}$; 
\item $\omega^2+\omega+1=0$. 
\end{enumerate*}



\item 已知复数$z$满足不等式$z\cdot\overline{z}+\mi{z}-\mi\overline{z}\leq 0$, 求$\abs{z}$的取值范围. 
\begin{jieda}
    $[0,2]$
\end{jieda}

\item 已知$z_1$,$ z_2$均为非零复数, 且$\abs{z_1+z_2}=\abs{z_1-z_2}$, 求证: $\left(\dfrac{z_1}{z_2}\right)^2$为负实数. 


\item (选做)已知$z^2+2z+3=\abs{z}^2$, 求复数$z$. 
\begin{jieda}
    $z=-\dfrac{3}{2}$或$-1+\frac{\sqrt{2}}{2}$或$-1-\frac{\sqrt{2}}{2}$
\end{jieda}
\end{enumerate}


\item 已知 $\left| z\right|  = 1$ ,且 ${z}^{2} + {2z} + \frac{1}{z}$ 是负实数,则复数 $z =$

\section{复数的几何意义}
\subsection{复数的几何意义（1）}
\zsd 复数的坐标表示与向量表示 \\
A组：
\begin{enumerate}
    \item 在复平面上,复数 $ \dfrac{5}{3 - {4i}} $ 所对应的点到原点的距离是\blkda{1}
    \begin{jieda}
        $ \because \dfrac{5}{3 - {4i}} = \dfrac{5\left( {3 + {4i}}\right) }{\left( {3 - {4i}}\right) \left( {3 + {4i}}\right) } = \dfrac{3}{5} + \dfrac{4}{5}i $ ,
        $ \therefore $ 复数 $ \frac{5}{3 - {4i}} $ 所对应的点到原点的距离是 $ \sqrt{{\left( \dfrac{3}{5}\right) }^{2} + {\left( \dfrac{4}{5}\right) }^{2}} = 1 $ .
        故答案为：1 .
    \end{jieda} 
    \item 若复数$z$满足$2z+\overline{z}=3-\mi$其中$\mi$为虚数单位, $\overline{z}$为$z$的共轭复数, 则$z$在复平面内对应的点位于第\blkda{四} 象限.
    \item 设$z=(a-2\mi)(1+3\mi)(a\zr)$, 其中$\mi$为虚数单位. 若$z$是纯虚数, 则$z$在复平面上对应点的坐标为\blkda{$(0,-20)$}. 
    \item 在复平面上, 已知$A(2,1)$, $B(0,2)$, $C(-2,1)$三点, $O$为原点. 给出下面结论: 
    \begin{dingautolist}{172}
    \item 直线$OC$与直线$BA$平行;
    \item $\vv{AB}+\vv{BC}=\vv{CA}$;
    \item $\vv{OA}+\vv{OC}=\vv{OB}$;
    \item $\vv{AC}=\vv{OB}-2\vv{OA}$. 
    \end{dingautolist}
    其中所有正确的结论的序号为\blkda{\ding{172}\ding{174}\ding{175}}. 
    \item 复数$z=\dfrac{3+4\mi}{1-2\mi}$在复平面内的对应点位于复平面的第\blkda{二}象限.
\end{enumerate}
B组：
\begin{enumerate}
    \item 已知复数 $ {z}_{1} = 5 - {2022i},{z}_{2} = {2017} + {2ai}\left( {a \in  \mathbf{R}}\right) $ ,若 $ {z}_{1} + {z}_{2} $ 所对应的点在实轴上,则 $ a = $
    \blkda{1011}.
    \begin{jieda}
        $ \because {z}_{1} + {z}_{2} $ 所对应的点在实轴上,
        $ \therefore {z}_{1} + {z}_{2} = 5 - {2022i} + {2017} + {2ai} = {2022} + \left( {{2a} - {2022}}\right) i $ ,
        $ \because {z}_{1} + {z}_{2} $ 所对应的点在实轴上,
        $ \therefore {2a} - {2022} = 0 $ ,解得 $ a = {1011} $ .
        故答案为：1011.
    \end{jieda}
    \item 已知复平面上有点 $ C\left( {2,4}\right) $ 和点 $ D $ ,使得向量 $ \overrightarrow{CD} $ 所对应的复数是 $ 1 + i $ ,则点 $ D $ 的坐标为\blkda{$(3, 5)$}
    \begin{jieda}
        解: 向量 $ \overrightarrow{CD} $ 所对应的复数是 $ 1 + i $ ,
        则 $ \overrightarrow{CD} = \left( {1,1}\right) $ ,
        设 $ D $ 点的坐标为$(x,y)$,
        则 $ x - 2 = 1,y - 4 = 1 $ ,解得 $ x = 3,y = 5 $ ,
        故点 $ D $ 的坐标为$(3, 5)$.
        故答案为：$(3, 5)$.
    \end{jieda}
    \item 已知复数$z_1=-1+2\mi$, $z_2=1-\mi$, $z_3=3-2\mi$, 它们所对应的点分别为$A,B,C$. 若$\vv{OC}=x\vv{OA}+y\vv{OB}$, 则$x+y$的值是\blkda{5}. 
    \item 在复平面上, 一个正方形的三个顶点对应的复数分别是$0$, $1+2\mi$, $-2+\mi$, 则该正方形的第四个顶点对应的复数是\blkda{$-1+3\mi$}. 
    \item 在复平面上, 已知正方形$OABC$(按逆时针方向, $O$表示原点)中的一个顶点$B$对应的复数为$1+2\mi$, 则$\vv{BC}$对应的复数$z=$\blkda{$-\dfrac{3}{2}-\dfrac{1}{2}\mi$}. 
    

    \item 已知复平面上有点 $ A $ 和点 $ B $ ,向量 $ \overrightarrow{OA} $ 与向量 $ \overrightarrow{AB} $ 所对应的复数分别为 $ - 1 - {2i} $ 与 $ 4 - i $ ,则点$ B $ 的坐标为\blkda{$\left( {3, - 3}\right)$}.
    \begin{jieda}
        $ \because \overrightarrow{OB} = \overrightarrow{OA} + \overrightarrow{AB} $ ,
        $ \therefore \overrightarrow{OB} $ 对应的复数为 $ - 1 - {2i} + 4 - i = 3 - {3i} $ ,
        故点 $ B $ 的坐标为$\left( {3, - 3}\right)$,
        故答案为:$\left( {3, - 3}\right)$.
    \end{jieda}
    \item 已知复平面上平行四边形 $ {ABCD} $ 的顶点 $ A\left( {-2, - 1}\right) ,B\left( {7,3}\right) \text{、}C\left( {{12},9}\right) \text{、}D\left( {x,y}\right) $ 按逆时针方向排列,则向量 $ \overrightarrow{AD} $ 所对应的复数为\blkda{$ 5 + {6i} $}.
    \begin{jieda}
        根据题意,复平面上平行四边形 $ {ABCD} $ 的顶点 $ A\left( {-2, - 1}\right) ,B\left( {7,3}\right) \text{、}C\left( {{12},9}\right) $ 、$ D\left( {x,y}\right) $ 按逆时针方向排列,
        则有 $ \overrightarrow{AB} = \overrightarrow{DC} $ ,
        而 $ \overrightarrow{AB} = \left( {9,4}\right) ,\overrightarrow{DC} = \left( {{12} - x,9 - y}\right) $ ,
        则有 $ \left( {9,4}\right)  = \left( {{12} - x,9 - y}\right) $ ,解可得 $ \left\{  \begin{array}{l} x = 3 \\  y = 5 \end{array}\right. $ ,
        故 $ D $ 的坐标为(3,5),则向量 $ \overrightarrow{AD} = \left( {5,6}\right) $ .
        所以对应的复数 $ Z = 5 + {6i} $ .
        故答案为: $ 5 + {6i} $ .
    \end{jieda}
    
    \item 已知等腰梯形$OABC$的顶点$A$、$B$在复平面上对应的复数分别为$1+2\mi$、$-2+6\mi$, 且$OA\px BC$, 则顶点$C$所对应的复数$z$为\blkda{$-5$}. 
    
    \item 若复数$z_1=4-3\mi$, $z_2=4+3\mi$(其中$\mi$为虚数单位)所对应的向量分别为$\vv{OZ_1}$与$\vv{OZ_2}$, 则$\triangle{OZ_1 Z_2}$的周长为\blkda{16}. 
    
    \item 已知复数$z$满足$\abs{z}=1$, 且$z^5+z=1$, 则复数$z=$\blkda{$\dfrac{1}{2}\pm\dfrac{\sqrt{3}}{2}\mi$}.
    \begin{jieda}
        \begin{enumerate}
            \item 方法1 \\
            $z^5 = 1-z$，两侧同时取模得到$|z^5| = |1-z|$，即$|1-z| = 1$，根据复数的几何意义$|1-z| = 1$对应圆心为$(1, 0)$的单位圆，$|z| = 1$对应圆心为$(0, 0)$的单位圆，两者交点为$z$所对应的点，因此$\dfrac{1}{2}\pm\dfrac{\sqrt{3}}{2}\mi$
            \item 方法2 \\
            设$z = cos\theta + i sin\theta$，则$z^5 = cos5\theta+isin5\theta$，因此$\left\{\begin{array}{l}
                 cos\theta + cos5\theta = 1 \\
                 sin\theta = -sin5\theta
            \end{array}\right.$，下面的三角方程解得$\theta = \dfrac{k\pi}{3}, k \in \mathbf{Z}$或$\theta = \dfrac{2k+1}{4}\pi, k \in \mathbf{Z}$，代入上面的三角方程检验可知$\theta = 2k\pi + \dfrac{\pi}{3}$或$2k\pi - \dfrac{\pi}{3}$，即$\dfrac{1}{2}\pm\dfrac{\sqrt{3}}{2}\mi$
        \end{enumerate}
    \end{jieda}
    
    \item 已知复数$z=3+4\mi$所对应的向量为$\vv{OZ}$, 把$\vv{OZ}$依逆时针旋转$\theta$得到一个新向量为$\vv{OZ_1}$. 若$\vv{OZ_1}$对应一个纯虚数, 
    当$\theta$取最小正角时, 这个纯虚数是\blkda{$5\mi$}. 
    
    \item 复数$6+5\mi$与$-3+4\mi$分别表示向量$\vv{OA}$, $\vv{OB}$, 则表示向量$\vv{6a}$的复数在复平面内对应的点位于\blkda{第一象限}

    \item 在复平面上, 复数$z=3-2\mi$对应的点关于直线$y=-x$对称的点所对应的复数为\blkda{$2-3\mi$}
    \item 实数$m$分别取什么数值时, 复数$z=\left(m^2+m-2\right)+(m-1)\mi$满足下列条件: 
    \begin{enumerate}[(1)]
    \item 纯虚数; 
    \item 对应的点在第一象限内. 
    \end{enumerate}
    \begin{jieda}
        \begin{enumerate}
            \item $\left\{\begin{array}{l}
                 m^2+m-2 = 0 \\
                 m - 1 \neq 0
            \end{array}\right.$，解得$m = -2$
            \item $\left\{\begin{array}{l}
                 m^2+m-2 > 0 \\
                 m - 1 > 0
            \end{array}\right.$，解得$m = (1, +\infty)$
        \end{enumerate}
    \end{jieda}
    
    
    \item 已知模为3的复数$z$所对应的向量为$\vv{OZ}$, 复数$z_1=3+2\mi$所对应的向量为$\vv{OZ_1}$, 求使得$\vv{OZ}\bot\vv{OZ_1}$成立的复数$z$. 
    \begin{enumerate}
        \item $\vv{OZ}\bot\vv{OZ_1}$即$\vv{OZ}\cdot \vv{OZ_1} = 0$，$\vv{OZ_1} = (3, 2)$，设$\vv{OZ} = (a, b)$，故$3a+2b = 0$，又因为$|z| = 3$，即$|\vv{OZ}| = 3$，因此$a^2+b^2 = 9$，消去$b$，得到$a^2 + \dfrac{9a^2}{4} = 9$，因此$\left \{\begin{array}{l}
             a = \dfrac{6\sqrt{13}}{13} \\
             b = -\dfrac{9\sqrt{13}}{13}
        \end{array} \right.$或$\left \{\begin{array}{l}
             a = -\dfrac{6\sqrt{13}}{13} \\
             b = \dfrac{9\sqrt{13}}{13}
        \end{array} \right.$，故$z = \dfrac{6\sqrt{13}}{13} - \dfrac{9\sqrt{13}}{13}i$或$z = -\dfrac{6\sqrt{13}}{13} + \dfrac{9\sqrt{13}}{13}i$
    \end{enumerate}
    
    \item 已知复数$z$的实部大于零, 且满足$z=\sqrt{2}\left(\cos\theta+\mi\sin\theta\right)\left(\theta\zr\right)$, $z^2$的虚部为2. 
    \begin{enumerate}[(1)]
    \item 求复数$z$; 
    \item 设$z$、$z^2$、$z-z^2$在复平面上的对应点分别为$A,B,C$, 求$\vv{AB}\cdot\vv{AC}$的值. 
    \end{enumerate}
    \begin{jieda}
        \begin{enumerate}
            \item $z^2 = 2(cos2\theta + i sin2\theta)$，因为$z^2$虚部为$2$，所以$sin2\theta = 1$，$\theta = k\pi + \dfrac{\pi}{4}, k \in \mathbf{Z}$，又因为$cos\theta > 0$，所以$\theta = 2k\pi + \dfrac{\pi}{4}, k \in \mathbf{Z}$，故$z = 1+i$
            \item $A(1, 1), B(0, 2), C(1, -1)$，故$\vv{AB} = (-1, 1), \vv{AC} = (0, -2)$，故$\vv{AB} \cdot \vv{AC} = -2$
        \end{enumerate}
    \end{jieda}
\end{enumerate}
C组：
\begin{enumerate}
    \item 如果 $ y = f\left( x\right) $ 的定义域为 $ R $ ,对于定义域内的任意 $ x $ ,存在实数 $ a $ 使得 $ f\left( {x + a}\right)  = f\left( {-x}\right) $
成立,则称此函数具有 “ $ P\left( a\right) $ 性质”. 给出下列命题：
\begin{dingautolist}{172}
    \item 函数 $ y = \sin x $ 具有“ $ P\left( a\right) $ 性质”；
    \item 若奇函数 $ y = f\left( x\right) $ 具有“ $ P\left( 2\right) $ 性质”,且 $ f\left( 1\right)  = 1 $ ,则 $ f\left( {2015}\right)  = 1 $ ；
    \item 若函数 $ y = f\left( x\right) $ 具有“ $ P\left( 4\right) $ 性质”,图象关于点(1,0)成中心对称,且在(-1,0)上单调递减,
    则 $ y = f\left( x\right) $ 在(-2, - 1)上单调递减,在(1,2)上单调递增；
    \item 若不恒为零的函数 $ y = f\left( x\right) $ 同时具有“ $ P\left( 0\right) $ 性质”和“ $ P\left( 3\right) $ 性质”,且函数 $ y = g\left( x\right) $ 对 $ \forall {x}_{1} $ ,
    $ {x}_{2} \in  \mathbf{R} $ ,都有 $ \left| {f\left( {x}_{1}\right)  - f\left( {x}_{2}\right) }\right|  \geq  \left| {g\left( {x}_{1}\right)  - g\left( {x}_{2}\right) }\right| $ 成立,则函数 $ y = g\left( x\right) $ 是周期函数.
\end{dingautolist}
其中正确的是\blkda{\ding{172}\ding{174}\ding{175}}(写出所有正确命题的编号).
\begin{jieda}
    \begin{dingautolist}{172}
        \item $ \because \sin \left( {x + \pi }\right)  =  - \sin x = \sin \left( {-x}\right) $ .
        $ \therefore $ 函数 $ y = \sin x $ 具有“ $ P\left( \pi \right) $ 性质”；
        即存在 $ a = \pi $ ,使函数 $ y = \sin x $ 具有“ $ P\left( a\right) $ 性质”； $ \therefore $ ，故正确
        \item  $ \because $ 若奇函数 $ y = f\left( x\right) $ 具有“ $ P\left( 2\right) $ 性质”,
        $ \therefore f\left( {x + 2}\right)  = f\left( {-x}\right)  =  - f\left( x\right) $ ,
        $ \therefore f\left( {x + 4}\right)  =  - f\left( {x + 2}\right)  = f\left( x\right) $ ,
        则函数 $ f\left( x\right) $ 的周期为 4 ,
        $ \because f\left( 1\right)  = 1,f\left( {2015}\right)  = f\left( 3\right)  = f\left( {-1}\right)  =  - f\left( 1\right)  =  - 1 $ ,
        故不正确.
        \item $ \because $ 若函数 $ y = f\left( x\right) $ 具有“ $ P\left( 4\right) $ 性质”,
        $ \therefore f\left( {x + 4}\right)  = f\left( {-x}\right) $ ,
        $ \therefore f\left( x\right) $ 关于 $ x = 2 $ 对称,
        即 $ f\left( {2 - x}\right)  = f\left( {2 + x}\right) $ .
        $ \because $ 图象关于点(1,0)成中心对称,
        $ \therefore f\left( {2 - x}\right)  =  - f\left( x\right) $ ,
        即 $ f\left( {2 + x}\right)  =  - f\left( {-x}\right) $ .
        $ \therefore $ 得出: $ f\left( x\right)  = f\left( {-x}\right) $ ,
        $ f\left( x\right) $ 为偶函数,
        又 $ f\left( x\right) $ 在(-1,0)上单调递减.
        $ \therefore f\left( x\right) $ 在(0,1)上单调增,又图象关于点(1,0)成中心对称,
        $ \therefore f\left( x\right) $ 在(1,2)上单调递增,
        $ \because f\left( x\right) $ 为偶函数, $ \therefore f\left( x\right) $ 在(-2, - 1)上单调递减,故正确.
        \item $ \because $ 函数 $ y = f\left( x\right) $ 同时具有“ $ P\left( 0\right) $ 性质”和“ $ P\left( 3\right) $ 性质”,
        $ \therefore f\left( x\right)  = f\left( {-x}\right) ,f\left( {x + 3}\right)  = f\left( {-x}\right)  = f\left( x\right) $ .
        $ \therefore f\left( x\right) $ 为偶函数,且周期为 3 ,
        $ \because $ 函数 $ y = g\left( x\right) $ 对 $ \forall {x}_{1},{x}_{2} \in  \mathbf{R} $ ,都有 $ \left| {f\left( {x}_{1}\right)  - f\left( {x}_{2}\right) }\right|  \geq  \left| {g\left( {x}_{1}\right)  - g\left( {x}_{2}\right) }\right| $ 成立,
        令 $ {x}_{1} = x + 3,{x}_{2} = x $ ,
        则 $ \left| {f\left( {x}_{1}\right)  - f\left( {x}_{2}\right) }\right|  = \left| {f\left( {x + 3}\right)  - f\left( x\right) }\right|  = 0 $ ,
        $ \therefore \left| {g\left( {x + 3}\right)  - g\left( x\right) }\right|  \leq  0 $ ,
        故 $ g\left( {x + 3}\right)  = g\left( x\right) $ ,因为 $ x $ 具有任意性,
        $ \therefore g\left( x\right) $ 是以 3 为周期的函数，即函数 $ y = g\left( x\right) $ 是周期函数,故正确.
    \end{dingautolist}
\end{jieda}
\item 已知函数 $ f\left( x\right)  = a\sin x + b\cos x + c\sin x\cos x + 1\left( {a,b,c \in  \mathrm{R}}\right) $ .
\begin{enumerate}
    \item 当 $ a = b = 0,c = 1 $ 时,求函数 $ y = f\left( x\right) $ 的单调增区间；
    \item 当 $ a = 1,c = 0 $ 时,设 $ g\left( x\right)  = f\left( x\right)  - 1 $ ,且函数 $ g\left( x\right) $ 的图像关于直线 $ x = \dfrac{\pi }{6} $ 对称,将函数 $ y = g\left( x\right) $ 的图像向右平移 $ \dfrac{\pi }{6} $ 个单位,得到函数 $ y = h\left( x\right) $ ,求解不等式 $ h\left( x\right)  \geq  1 $ ;
    \item 当 $ a = b = 1,c = 0 $ 时,若实数 $ m,n,p $ 使得 $ {mf}\left( x\right)  + {nf}\left( {x - p}\right)  = 1 $ 对任意实数 $ x $ 恒成立,求实数 $ m,n,p $ 的值.
\end{enumerate}
\begin{jieda}
    \begin{enumerate}
        \item 当 $ a = b = 0,c = 1 $ 时,可得函数 $ f\left( x\right)  = \sin x\cos x + 1 = \dfrac{1}{2}\sin {2x} + 1 $ ,
        令 $ {2k\pi } - \dfrac{\pi }{2} \leq  {2x} \leq  {2k\pi } + \dfrac{\pi }{2},k \in  \mathbf{Z} $ ,所以单调增区间为 $ \left\lbrack  {{k\pi } - \dfrac{\pi }{4},{k\pi } + \dfrac{\pi }{4}}\right\rbrack  ,k \in  \mathbf{Z} $ ;
        \item 当 $ a = 1,c = 0 $ 时,可得 $ g\left( x\right)  = f\left( x\right)  - 1 = \sin x + b\cos x = \sqrt{1 + {b}^{2}}\sin \left( {x + \theta }\right) $ ,其中 $ \tan \theta  = b $ ,
        因为 $ g\left( x\right) $ 关于直线 $ x = \dfrac{\pi }{6} $ 对称,可得 $ g{\left( x\right) }_{\max } = \left| {g\left( \dfrac{\pi }{6}\right) }\right|  = \sqrt{1 + {b}^{2}} $ ,即 $ \left| {\dfrac{1}{2} + \dfrac{\sqrt{3}}{2}b}\right|  = \sqrt{1 + {b}^{2}} $ ,解得 $ b = \sqrt{3} $ ,
        所以 $ g\left( x\right)  = \sin x + \sqrt{3}\cos x = 2\sin \left( {x + \dfrac{\pi }{3}}\right) $ ,将函数 $ y = g\left( x\right) $ 的图像向右平移 $ \dfrac{\pi }{6} $ 个单位,得到函数$ h\left( x\right)  = 2\sin \left( {x + \dfrac{\pi }{6}}\right) $ ,由 $ h\left( x\right)  \geq  1 $ ,即 $ \sin \left( {x + \dfrac{\pi }{6}}\right)  \geq  \dfrac{1}{2} $ ,则 $ \dfrac{\pi }{6} + {2k\pi } \leq  x + \dfrac{\pi }{6} \leq  \dfrac{5}{6}\pi  + {2k\pi },k \in  Z $
        解得 $ {2k\pi } \leq  x \leq  \dfrac{2}{3}\pi  + {2k\pi },k \in  Z $ ,所以不等式的解集为 $ \left\lbrack  {{2k\pi },\dfrac{2}{3}\pi  + {2k\pi }}\right\rbrack  \left( {k \in  Z}\right) $ 
        \item 当 $ a = 3,b = 2,c = 0 $ 时,则 $ f\left( x\right)  = \sin x + \cos x + 1 $ ,
        可得 $ f\left( x\right)  = \sqrt{2}\sin \left( {x + \dfrac{\pi }{4}}\right)  + 1 $ ,则 $ f\left( {x - p}\right)  = \sqrt{2}\sin \left( {x + \dfrac{\pi }{4} - p}\right)  + 1 $ ,
        于是 $ {mf}\left( x\right)  + {nf}\left( {x - p}\right)  = 1 $ ,
        可化为 $ \sqrt{2}m\sin \left( {x + \dfrac{\pi }{4}}\right)  + \sqrt{2}n\sin \left( {x + \dfrac{\pi }{4} - p}\right)  + m + n = 1 $ ,
        即 $ \sqrt{2}m\sin \left( {x + \dfrac{\pi }{4}}\right)  + \sqrt{2}n\sin \left( {x + \dfrac{\pi }{4}}\right) \cos p - \sqrt{2}n\sin p\cos \left( {x + \dfrac{\pi }{4}}\right)  + \left( {m + n - 1}\right)  = 0 $ .
        所以 $ \sqrt{2}\left( {m + n\cos p}\right) \sin \left( {x + \dfrac{\pi }{4}}\right)  - \sqrt{2}n\sin p\cos \left( {x + \dfrac{\pi }{4}}\right)  + \left( {m + n - 1}\right)  = 0 $ .
        由已知条件,上式对任意 $ x \in  \mathrm{R} $ 恒成立,故必有 $ \left\{  \begin{array}{l} m + n\cos p = 0\cdots \left( 1\right) \\  n\sin p = 0\cdots \left( 2\right) \\  m + n - 1 = 0\cdots \left( 3\right)  \end{array}\right. $ ,
        若 $ n = 0 $ ,则由(1)知 $ m = 0 $ ,显然不满足(3)式,故 $ n \neq  0 $ ,
        所以由 (2) 知 $ \sin p = 0 $ ,故 $ p = {2k\pi } + \pi $ 或 $ p = {2k\pi },k \in  \mathbf{Z} $ ,
        当 $ p = {2k\pi } $ 时, $ \cos p = 1 $ ,则(1)、(3)两式矛盾,
        故 $ p = {2k\pi } + \pi ,k \in  Z,\cos p =  - 1 $ 由(1)、(3)知 $ m = n = \dfrac{1}{2} $ ,
        所以 $ p = {2k\pi } + \pi ,k \in  \mathbf{Z},m = n = \dfrac{1}{2} $ ;
    \end{enumerate}
\end{jieda}
\end{enumerate}

\clearpage
\subsection{复数的几何意义（2）}
\zsd 复数的模与几何意义 \\
\begin{enumerate}
    \item 已知复数 $ z $ 满足 $ z \cdot  {\left( 2 - i\right) }^{2} = 2 + i(i $ 为虚数单位),则 $ \left| z\right| = $ \blkda{$\dfrac{\sqrt{5}}{5}$}
    \begin{jieda}
        $ z = \dfrac{2 + i}{{\left( 2 - i\right) }^{2}} = \dfrac{2 + i}{3 - {4i}} = \dfrac{\left( {2 + i}\right) \left( {3 + {4i}}\right) }{\left( {3 - {4i}}\right) \left( {3 + {4i}}\right) } = \dfrac{2}{25} + \dfrac{11}{25}i $ ,所以 $ \left| z\right|  = \sqrt{{\left( \dfrac{2}{25}\right) }^{2} + {\left( \dfrac{11}{25}\right) }^{2}} = \dfrac{\sqrt{5}}{5} $ .
    \end{jieda}
    \item 若复数$z=m^2-4m-5+\left(m+\dfrac{1}{m}-2\right)\mi$在复平面内对应的点在第二象限, 则实数$m$的取值范围是\blkda[3]{$(0,1)\cup(1,5)$}. 

    \item 若复数$z=\dfrac{m-1}{3}-(m-2)\mi(m\zr)$, 它在复平面上对应的点为$Z$, 则复平面上的点$(1,2)$到点$Z$之间的最短距离是\blkda{$\dfrac{2\sqrt{10}}{5}$}. 
    \item 已知复数 $ z $ 且 $ \left| z\right|  = 1 $ ,则 $ \left| {z - 2 - {2i}}\right| $ 的最小值是\blkda{$2\sqrt{2} - 1$}
    \begin{jieda}
        $ \therefore \left| z\right|  = 1 $ ,则 $ z $ 对应的点 $ Z $ 在以原点为圆心,以 1 为半径的圆上,如图所示，
        \begin{center}
    \includegraphics[width=0.3\textwidth]{bodyxb//src/fushujiheyiyi1.png}
\end{center}\\
        可得$\left| {z - 2 - {2i}}\right| $ 的最小值是复数 $ 2 + {2i} $ 对应的点(2,2)到原点的距离减去半径 1 ,
        即 $ \sqrt{{2}^{2} + {2}^{2}} - 1 = 2\sqrt{2} - 1 $ .
    \end{jieda}
    
    \item 设 $ a \in  \mathbf{R},i $ 为虚数单位. 若对于任意 $ \theta  \in  \mathbf{R} $ ,复数 $ z = \left( {a - \cos \theta }\right)  + \left( {a - 1 - \sin \theta }\right) i $ 的模始终不大于 2,则 $ a $ 的取值范围是\blkda{$\left\lbrack  {0,1}\right\rbrack$}.
    \begin{jieda}
        $ \left| z\right|  = \left| {a + \left( {a - 1}\right) i - \left( {\cos \theta  + i\sin \theta }\right) }\right|  \leq  \left| {a + \left( {a - 1}\right) i}\right|  + \left| {\cos \theta  + i\sin \theta }\right|  = \sqrt{{a}^{2} + {\left( a - 1\right) }^{2}} + 1 \leq  2 $ ,
        $ \therefore {a}^{2} + {\left( a - 1\right) }^{2} \leq  1,\;\therefore 0 \leq  a \leq  1 $ .
        故答案为: $ \left\lbrack  {0,1}\right\rbrack $ .
    \end{jieda}
    \item 已知$z\zc$, 且$\abs{z+3+4\mi}\leq 2$, 则$\abs{z}$的最大值为\blkda{7}. 


    \item 虚数$z=(x-2)+y\mi(x,y\zr)$, 且$\abs{z}=1$, 则$\dfrac{y}{x}$的取值范围为\blkda[3]{$[-\frac{\sqrt{3}}{3},0)\cup(0,\frac{\sqrt{3}}{3}]$}. 
    
    
    \item 复数$z_1 =1+\mi\cos\theta$, $z_2=\sin\theta-\mi$, 则$\abs{z_1-z_2}$的最大值为\blkda{$1+\sqrt{2}$}. 
    
    \item 已知复数$z$满足$\abs{z}=1$, 则$\abs{z-2\mi}$(其中$\mi$是虚数单位)的最小值为\blkda{1}. 


    \item 若集合$M=\{z\mid\abs{z+1}=1, z\zc\}$, 集合$N=\{z\mid\abs{z-2\mi}=\abs{z}, z\zc\}$, 则$M\cap N=$\blkda{$\{-1+\mi\}$}. 
    
    \item 若复数$z$满足$\abs{z+\mi}-\abs{z}=1$, 则$\abs{z+1+\mi}$的最小值为\blkda{$\sqrt{2}$}. 
    
    \item 如果复数$z$满足$\abs{z+\mi}+\abs{z-\mi}=2$, 那么$\abs{z+\mi+1}$的最小值为\blkda{1}.

    \item 若$\mi$为虚数单位, 复数$z$满足$1\leq\abs{z+1+\mi}\leq\sqrt{2}$, 则$\abs{z-1-\mi}$的最大值为\blkda{$3\sqrt{2}$}. 

    
    \item 如图, 复平面中的阴影部分(含边界)对应的复数集合是\dotfill(\kh{D})
    \begin{center}
    \includegraphics[width=0.3\textwidth]{bodyxb//src/fushudejiheyiyi2.png}
\end{center}
    \xx
    {$\left\{z\mid\abs{z}=1, \Re{z}\geq\dfrac{1}{2}, z\zc\right\}$}
    {$\left\{z\mid\abs{z}\leq1, \Re{z}\geq\dfrac{1}{2}, z\zc\right\}$}
    {$\left\{z\mid\abs{z}=1, \Im{z}\geq\dfrac{1}{2}, z\zc\right\}$}
    {$\left\{z\mid\abs{z}\leq1, \Im{z}\geq\dfrac{1}{2}, z\zc\right\}$}
    \item 复数$z$满足$z\cdot\mi=\abs{\dfrac{1}{2}+\dfrac{\sqrt{3}}{2}\mi}$, 则复数$z$在复平面内对应的点的坐标为\dotfill(\kh{D})
    \xx{$(1,0)$}{$(0,1)$}{$(-1,0)$}{$(0,-1)$}
    \item 已知 $ i $ 为虚数单位,下列说法中错误的是\dotfill{(\kh{D})}
    \xx{复数 $ {z}_{1} $ 对应的向量为 $ \vv{O{Z}_{1}} $ ,复数 $ {z}_{2} $ 对应的向量为 $ \vv{O{Z}_{2}} $ ,若 $ \left| {{z}_{1} + {z}_{2}}\right|  = \left| {{z}_{1} - {z}_{2}}\right| $ ,则 $\vv{O{Z}_{1}} \bot  \vv{O{Z}_{2}} $}
    {互为共轭复数的两个复数的模相等,且 $ {\left| \bar{z}\right| }^{2} = {\left| z\right| }^{2} = z \cdot  \bar{z} $}
    {复数的模实质上就是复平面内复数对应的点到原点的距离,也就是复数对应的向量的模}
    {若复数 $ z $ 满足 $ \left| {z - i}\right|  = \sqrt{5} $ ,则复数 $ z $ 对应的点在以(1,0)为圆心, $ \sqrt{5} $ 为半径的圆上}
    \begin{jieda}
        \begin{enumerate}
            \item A的几何意义是以$\vv{OZ_1}, \vv{OZ_2}$为临边的平行四边形的两对角线相等，因此其为矩形，故正确
            \item 根据模的几何意义是复平面上的点到原点的距离可知B正确
            \item 同上
            \item 圆心应为$(0, 1)$
        \end{enumerate}
    \end{jieda}
    \item 设$z_1, z_2 \in \mathbf{C}$，已知$|z_1| = 3, |z_2| = 5, |z_1 - z_2| = \sqrt{10}$，则$|z_1+z_2| = $\blkda{$\sqrt{58}$}
    \begin{jieda}
        \begin{enumerate}
            \item 代数方法 \\
            $|z_1| = 3$即$z_1 \overline{z_1} = 9$，$|z_2| = 5$即$z_2 \overline{z_2} = 25$，$|z_1 - z_2| = \sqrt{10}$即$(z_1 - z_2)(\overline{z_1} - \overline{z_2}) = 10$ \\
            故$z_1\overline{z_2} + \overline{z_1}z_2 = 24$，故$(z_1 + z_2)(\overline{z_1} + \overline{z_2}) = 58, |z_1+z_2| = \sqrt{58}$
            \item 几何方法 \\
            设$z_1, z_2$在复平面上对应的点为$Z_1, Z_2$，则$|\vv{OZ_1}| = 3, |\vv{OZ_2}|=5, |\vv{Z_2Z_1}| = \sqrt{10}$，设$\vv{OZ} = \vv{OZ_1} + \vv{OZ_2}$，则$|z_1 + z_2| = |\vv{OZ}|$，根据几何关系有$|\vv{OZ}|^2 + |\vv{Z_2Z_1}|^2 = 2(|\vv{OZ_1}|^2+|\vv{OZ_2}|^2)$，因此$|\vv{OZ}|^2 = 2\cdot(9+25) - 10 = 58$，因此$|z_1+z_2| = \sqrt{58}$
        \end{enumerate}
        
        \item 已知复数$z$满足$z=1+\mi$, 复数$z$的共轭复数为$\overline{z}$. 
        \begin{enumerate}[(1)]
        \item 求$\dfrac{\overline{z}}{z}+z^4$; 
        \item 若复数$z_1$满足$\abs{z_1-z}=1$, 求$\abs{z_1}$的最小值和最大值. 
        \end{enumerate}
        \begin{jieda}
            \begin{enumerate}
                \item $\overline{z} = 1-i$，故$\dfrac{\overline{z}}{z} = \dfrac{1-i}{1+i} = -i$，$z^4 = (\sqrt{2})^4[cos\pi + i sin\pi] = -4$，故$\dfrac{\overline{z}}{z}+z^4 = -4-i$
                \item 根据几何意义可知$z_1$在复平面上对应的点在以$(1, 1)$为圆心，1为半径的圆上，因此$|z_1| \in [\sqrt{2}-1, \sqrt{2}+1]$
            \end{enumerate}
        \end{jieda}
        
        
        \item 如图, 设复数$z=(2a-1)+\sqrt{2}a\mi, a\zr$在复平面内的对应点落在如图所示的阴影部分内(包括边界), 求实数$a$的取值范围. 
        \begin{center}
    \includegraphics[width=0.3\textwidth]{bodyxb//src/fushudejiheyiyi3.png}
\end{center}
        \begin{jieda}
            $|z| \geq 1$，$0 \leq Imz \leq 1$，$0 \leq Rez \leq 1$，即$\left\{\begin{array}{l}
                 4a^2-4a+1+2a^2 \geq 1 \\
                 0 \leq 2a-1 \leq 1 \\
                 0 \leq \sqrt{2}a \leq 1
            \end{array}\right.$，即$\left\{\begin{array}{l}
                 6a^2 - 4a \geq 0 \\
                 \dfrac{1}{2} \leq a \leq 1 \\
                 0 \leq a \leq \dfrac{\sqrt{2}}{2}
            \end{array}\right.$，故$a \in \left[\dfrac{2}{3}, \dfrac{\sqrt{2}}{2}\right]$
        \end{jieda}
        
        
        \item (选做)已知复平面上点集$S=\left\{z\mid\abs{z}^2-2\mi z+2a(1+\mi)=0\right\}, a\geq0$. 
        \begin{enumerate}[(1)]
        \item 若$S\neq\varnothing$, 求$a$的范围; 
        \item 当$S\neq\varnothing$时, 求$\abs{z-2}$的范围. 
        \end{enumerate}
        \begin{jieda}
            \begin{enumerate}
                \item 设$z = x+yi$，故$\left \{
                \begin{array}{l}
                     x^2+y^2+2y+2a=0  \\
                     -2x+2a=0
                \end{array}
                 \right.$，下式代入上式得到$a^2+2a+y^2+2y = 0$，即$(a+1)^2+(y+1)^2 = 2$，根据几何意义可知点$(a, y)$在以$(-1, -1)$为圆心，$\sqrt{2}$为半径的圆上，因为$a \geq 0$，所以$a \in [0, \sqrt{2}-1]$
                 \item 由上一问可知$z$在复平面内对应的点$(x, y)$在以$(-1, -1)$为圆心，$\sqrt{2}$为半径的圆在直线$x = 0$右侧圆弧上，故$z-2$在复平面上对应的点$(x-2, y)$在以$(-3, -1)$为圆心，$\sqrt{2}$为半径的圆在直线$x = -2$右侧圆弧上，故数形结合可知$|z-2|$的取值范围为$[\sqrt{10}-\sqrt{2}, 2\sqrt{2}]$
            \end{enumerate}
        \end{jieda}
    \end{jieda}
    \item (选做)设 $ 0 \leq  x < {2\pi } $ ,满足 $ \sin \left( {x + \frac{\pi }{6}}\right)  = \sin x + \sin \frac{\pi }{6} $ 的 $ x $ 的个数为\blkda{2}.
    \begin{jieda}
        由 $ \sin \left( {x + \dfrac{\pi }{6}}\right)  = \sin x + \sin \dfrac{\pi }{6} $ ,得 $ \sin \left( {x + \dfrac{\pi }{6}}\right)  - \sin x = 2\cos \left( {x + \dfrac{\pi }{12}}\right) \sin \dfrac{\pi }{12} = \sin \dfrac{\pi }{6} = \dfrac{1}{2} $
        即 $ x $ 在 $ \lbrack 0,{2\pi }) $ 内使 $ cos\left( {x + \dfrac{\pi }{12}}\right)  = \dfrac{1}{4\sin \dfrac{\pi \ }{12}} $ 成立 $ \theta $ 个数
        
        再根据 $ 0 < \sin \dfrac{\pi }{12} < \dfrac{1}{2} $ 知 $ 0 < \dfrac{1}{4\sin \dfrac{\pi }{12}} < \dfrac{1}{2} $ ,且 $ cos\dfrac{\pi }{12} = cos\left(\dfrac{25}{12}\pi\right)  > \dfrac{1}{2} $ , 故有 2 个
    \end{jieda}
    \item (选做) 设函数 $ f\left( x\right)  = a\left( {\cos {2x} + \cos x + 1}\right)  - \cos x - 1,g\left( x\right)  =  - $
    $ {2a}\sin {2x} + \left( {1 - a}\right) \sin x $ ,其中 $ a > 0 $ .
    \begin{enumerate}
        \item 当 $ a = 2 $ 时,求函数 $ f\left( x\right) $ 的值域
        \item 记 $ |f\left( x\right)| $ 的最大值为 $ M $ , \ding{172}求 $ M $ \ding{173}求证：
        $ \left| g(x)\right| \leq  {2M} $ .
    \end{enumerate}
    \begin{jieda}
        \begin{enumerate}
        \item $ a = 2 $ 时, $ f\left( x\right)  = 2\left( {\cos {2x} + \cos x + 1}\right)  - \cos x - 1 = 4{\cos }^{2}x + \cos x - 1$，令$t = cosx, t \in  \left\lbrack  {-1,1}\right\rbrack$，则$f(x) = 4{t}^{2} + t - 1$
        $ \therefore f\left( x\right)  \in  \left\lbrack  {-\dfrac{17}{16},4}\right\rbrack $
        \item 
        \ding{172} $ f\left( x\right)  = a\left( {\cos {2x} + \cos x + 1}\right)  - \cos x - 1 = {2a}{\cos }^{2}x + \left( {a - 1}\right) \cos x - 1 $\\
        令$t = cosx$，则$f(x) = g(t) = {2a}{t}^{2} + \left( {a - 1}\right) t - 1\;t \in  \left\lbrack  {-1,1}\right\rbrack $
        对称轴为 $ t = \dfrac{1 - a}{4a} $ ,且 $ \Delta  = {\left( a - 1\right) }^{2} + {8a} > 0 $
        \begin{enumerate}
            \item $ \dfrac{1 - a}{4a} <  - 1 $，即$ a <  - \dfrac{1}{3} $ 不符合题意
            \item $ \dfrac{1 - a}{4a} > 1 $，即$ 0 < a < \dfrac{1}{5} $ ,此时 $ \left| {g\left(-1\right) }\right|  = a < \left| {{3a} - 2}\right|  = \left| {g\left( 1\right) }\right| $ , $ M = 2 - {3a} $
            \item 
            $ - 1 \leq  \dfrac{1 - a}{4a} \leq  1 \Rightarrow  a \geq  \dfrac{1}{5} $ ,此时 $ \left| {g\left( {-1}\right) }\right|  = a $， $ \left| {g\left( 1\right) }\right|  = \left| {{3a} - 2}\right| $ , $ \left| {g\left( \dfrac{1 - a}{4a}\right) }\right|  = \dfrac{{a}^{2} + {6a} + 1}{8a} $ \\
            当 $ \dfrac{1}{5} \leq  a \leq  1 $ 时, $ \dfrac{{a}^{2} + {6a} + 1}{8a} - \left( {2 - {3a}}\right)  > 0 $ ,且 $ \dfrac{{a}^{2} + {6a} + 1}{8a} - a > 0 $ ,所以 $ M = \dfrac{{a}^{2} + {6a} + 1}{8a} $ \\
            当 $ a > 1 $ 时, $ a < {3a} - 2 $ 且 $ {3a} - 2 > \dfrac{{a}^{2} + {6a} + 1}{3a},\;M = {3a} - 2 $ .
        \end{enumerate}
        \ding{173} $ \left| {g\left( x\right) }\right|  = \left| {-{2a}\sin {2x} + \left( {1 - a}\right) \sin x}\right|  \leq  \left| {-{2a}\sin {2x}}\right|  + \left| {\left( {1 - a}\right) \sin x}\right|  \leq  {2a} + \left| {1 - a}\right| $
            $ = \left\{  \begin{array}{ll} a + 1 & 0 < a \leq  1 \\  {3a} - 1 & a > 1 \end{array}\right. $
            
            当 $ a > 1 $ 时, $ \left| {g\left( x\right) }\right|  - {2M} = {3a} - 1 - 2\left( {{3a} - 2}\right)  =  - {3a} + 3 < 0 $
            
            当$ \dfrac{1}{5} \leq  a \leq  1 $ 时 $ \left| {g\left( x\right) }\right|  - {2M} = a + 1 - \dfrac{{a}^{2} + {6a} + 1}{4a} = \dfrac{\left( {{3a} + 1}\right) \left( {a - 1}\right) }{4a} < 0 $
            
            当$ 0 < a < \dfrac{1}{5} $ 时 $ \left| {g\left( x\right) }\right|  - {2M} = a + 1 - 2\left( {2 - {3a}}\right)  = {7a} - 3 < 0 $ . 
            
            结论得证.
    \end{enumerate}
    \end{jieda}
\end{enumerate}

\section{复数集内的方程}
\subsection{复数集内的方程(1)}
\zsd 复数集内的方程 \\
A组：
\begin{enumerate}
    \item 若 $ {x}^{2} + x + 1 = 0 $ ,则$x = $\blkda[2]{$-\dfrac{1}{2}\pm \dfrac{\sqrt{3}}{2}i$}
    \item 若$1-\mi$是实系数方程$x^2+px+q=0$的根, 则$p\cdot q=$\blkda{$-4$}. 
    \item 若方程$x^2+ax+b=0$有一个根是$-1+\sqrt{3}\mi$, 则实数$a=$\blkda{2}, 实数$b=$\blkda{4}. 
    \item 定义: 若 $ {z}^{2} = a + b\mathrm{i}\left( {a,b \in  \mathbf{R}}\right) $ ,则称复数 $ z $ 是复数 $ a + b\mathrm{i} $ 的平方根. 根据定义, 复数 $ 9 - {40}\mathrm{i} $ 的平方根为\dotfill{(\kh{C})}
    \xx{$ 3 - 4\mathrm{i}, - 3 + 4\mathrm{i} $}{$ 4 + 3\mathrm{i},4 - 3\mathrm{i} $}{$ 5 - 4\mathrm{i}, - 5 + 4\mathrm{i} $}{$ 4 - 5\mathrm{i}, - 4 + 5\mathrm{i} $}
    \begin{jieda}
        设$z = a + b\mathrm{i}\left( {a,b \in  \mathbf{R}}\right) $满足$z^2 = 9-40i$，则$\left\{\begin{array}{l}
             a^2-b^2 = 9 \\
             2ab = -40
        \end{array}\right.$，解得$\left\{\begin{array}{l}
             a = 5 \\
             b = -4
        \end{array}\right.$或$\left\{\begin{array}{l}
             a = -5 \\
             b = 4
        \end{array}\right.$
    \end{jieda}
\end{enumerate}
B组：
\begin{enumerate}
    \item 若方程$x^2-ax+2=0(a\zr)$有虚根$z$, 则$\abs{z}$等于\blkda{$\sqrt{2}$}. 

    \item 若实系数方程$x^2+mx+m=0$有虚根, 则实数$m$的取值范围是\blkda{$(0,4)$}. 

    \item 若$x_1$, $x_2$为方程$x^2-x+7=0$的两个解, 则$\abs{x_1-x_2}^2=$\blkda{27}. 
    
    \item 设$z=\dfrac{1}{2}+\dfrac{\sqrt{3}}{2}\mi$, 那么$z+z^2+z^3+z^4+z^5+z^6=$\blkda{0}. 
    
    \item 已知方程$x^2-2\sqrt{2}x+m=0$的解为$x_1,x_2$, 且$\abs{x_1-x_2}=3$, 则实数$m$的值为\blkda[2]{$-\dfrac{1}{4}$或$\dfrac{17}{4}$}. 

    \item 对于实系数一元二次方程 $ a{x}^{2} + {bx} + c = 0 $ ,在复数范围内其解是 $ {x}_{1}\text{、}{x}_{2} $ ,下列结论中不正确的是\dotfill{(\kh{D})}
    \xx{若 $ {b}^{2} - {4ac} = 0 $ ,则 $ {x}_{1} = {x}_{2} \in  \mathbf{R} $}
    {若 $ {b}^{2} - {4ac} < 0 $ ,则 $ {x}_{1} \notin  \mathbf{R},{x}_{2} \notin  \mathbf{R} $ 且 $ x_1 = \overline{{x}_{2}} $}
    {一定有 $ {x}_{1} + {x}_{2} =  - \dfrac{b}{a},{x}_{1}{x}_{2} = \dfrac{c}{a} $}
    {一定有 $ {\left( {x}_{1} - {x}_{2}\right) }^{2} = \dfrac{\left| {b}^{2} - 4ac\right| }{{a}^{2}} $}
    \begin{jieda}
        D应为：一定有 $ {\left( {x}_{1} - {x}_{2}\right) }^{2} = \dfrac{{b}^{2} - 4ac }{{a}^{2}} $
    \end{jieda}


    \item 若$\omega=-\dfrac{1}{2}+\dfrac{\sqrt{3}}{2}\mi$, 则$(1+\omega)\cdot(1+\omega-\omega^2)=$\dotfill(\kh{D})
    \xx{$\omega^2$}{$-\omega$}{$-2\omega^2$}{$2\omega$}

    \item 若正整数$n$满足$\left(-\dfrac{1}{2}+\dfrac{\sqrt{3}}{2}\mi\right)^n+\mi^n=0$, 则$n$一定是\dotfill(\kh{B})
    \xx{3的倍数}{6的倍数}{奇数}{3的奇倍数}

    \item 在复数集内方程 $ {x}^{2} - \mathrm{i}x + \mathrm{i} - 1 = 0 $ 的解为\blkda[2]{1或$-1+i$}

    \begin{jieda}
        首先观察到$1$是一根，根据韦达定理可以得到另一根为$i-1$
    \end{jieda}

    \item 已知关于 $ x $ 的方程 $ {x}^{2} + \left( {k + 2\mathrm{i}}\right) x + 2 + k\mathrm{i} = 0 $ 有实数根,则实数 $ k $ 的值为\blkda{$\pm 2\sqrt{2}$}

    \begin{jieda}
        若$x$为实数，方程左侧的虚部$2x+k = 0$，故$x = -\dfrac{k}{2}$，因此方程左侧实部可化为$\dfrac{k^2}{4}+k\cdot\left(-\dfrac{k}{2}\right)+2 = 0$，解得$k = \pm 2\sqrt{2}$
    \end{jieda}

    \item 若 $ 2 - 3\mathrm{i} $ 是方程 $ {x}^{2} - {4x} + a = 0\left( {a \in  \mathrm{R}}\right) $ 的一个根,则其另外一个根是\blkda{$2+3i$}, $a = $ \blkda{13}

    \item 方程 $ {z}^{2} - 5\left| z\right|  + 6 = 0 $ 在复数集中解的个数为\blkda{6}
    \begin{jieda}
        ${z}^{2} - 5\left| z\right|  + 6 = 0 \Longleftrightarrow z^2 = 5|z| - 6$，因此$z$是实数或者纯虚数
        \begin{dingautolist}{172}
            \item 设$z = a, a \in \mathbf{R}$，则$a^2 = 5|a|-6$，即$\left|a + \dfrac{6}{a}\right| = 5$，有四个根：$\pm 2, \pm 3$
            \item 设$z = bi, b \in \mathbf{R}$，则$-b^2 = 5|b|-6$，即$\left|b - \dfrac{6}{b}\right| = 5 (b \in [-\sqrt{6}, 0) \cup(0, \sqrt{6}])$，有两个根：$\pm 1$
        \end{dingautolist}
    \end{jieda}

    \item 方程$4x^2+mx+1=0(m\zr)$, 则下列结论正确的是\dotfill(\kh{C})
    \xx{方程的两根互为共轭复数}
    {若方程的两根互为共轭复数, 则$m=0$}
    {若$x$为方程的一个虚根, 则$\overline{x}$也为方程的} 
    {若$m<0$, 则方程的两根一定为正数}

    \item 已知$z_1,z_2$是二次方程$z^2-2z+2=0$的两根, 且$z_1$在第一象限, 则$\dfrac{z_1}{z_2}=$\blkda{$\mi$}

    \item 已知复数 $ z = m + 2 + \left( {m - 2}\right) \mathrm{i}\left( {m \in  \mathbf{R}}\right) ,\bar{z} $ 为 $ z $ 的共轭复数,且 $ z + \bar{z} = 6 $ .
    \begin{enumerate}
        \item 求 $ m $ 的值；
        \item 若 $ z - 3\mathrm{i} $ 是关于 $ x $ 的实系数一元二次方程 $ {x}^{2} + {ax} + b = 0 $ 的一个根, 求该一元二次方程的另一复数根.
    \end{enumerate}
    \begin{jieda}
        \begin{enumerate}
            \item $z + \overline{z} = 6$即$Rez = 3$，因此$m = 1$，故$z = 3 - i$
            \item $z - 3i = 3 - 4i$为方程一根，因此另一个复数根为$3+4i$
        \end{enumerate}
    \end{jieda}

    \item 已知复数$\omega$满足$\omega-4=(3-2\omega)\mi$($\mi$为虚数单位), $z=\dfrac{5}{\omega}+\abs{\omega-2}$, 求一个以$z$为根的实系数一元二次方程. 
    \begin{jieda}
        $\omega$满足$\omega-4=(3-2\omega)\mi$整理得$\omega(1+2i) = 4+3i$，即$\omega = \dfrac{4+3i}{1+2i} = \dfrac{(4+3i)(1-2i)}{5} = 2-i$，因此$z = 3+i$，因此实系数一元二次方程的另一根为$3-i$，因此该方程可为$x^2 - 6x + 10 = 0$
    \end{jieda}

    \item 已知$\alpha$是非零复数$z$的一个立方根, 求证: $z$的另外两个立方根分别是$\alpha\omega$、$\alpha\overline{\omega}$, 其中$\omega=-\dfrac{1}{2}+\dfrac{\sqrt{3}}{2}\mi$. 
    \begin{jieda}
        $\omega^3 = 1, \overline{\omega}^3 = 1$，根据复数乘法的几何意义，结论显然.
    \end{jieda}
    
    \item 已知$x_1, x_2$是关于$x$的方程$x^2+2x+m=0(m \in \mathbf{R})$的两个根，求$|x_1| + |x_2|$的值.
        \begin{jieda}
            $\Delta = 4 - 4m$
            \begin{dingautolist}{172}
                \item $\Delta \geq 0$，即$m \leq 1$ \\
                此时$x_1, x_2$为实数， $x_1 = -1+\sqrt{1-m}, x_2 = -1-\sqrt{1-m}$，故$|x_1|+|x_2| = \left \{ \begin{array}{ll}
                     2\sqrt{1-m}, &  m < 0 \\
                     2, & m \in [0, 1]
                \end{array} \right.$
                \item $\Delta < 0$，即$m > 1$ \\
                此时$x_1, x_2$为虚数， $x_1 = -1+i\sqrt{m-1}, x_2 = -1-i\sqrt{m-1}$，故$|x_1|+|x_2| = 2\sqrt{m}$
            \end{dingautolist}
            综上$|x_1|+|x_2| = \left \{ \begin{array}{ll}
                     2\sqrt{1-m}, &  m < 0 \\
                     2, & m \in [0, 1] \\
                     2\sqrt{m}, & m > 1
                \end{array} \right.$
        \end{jieda}
    

    \item 已知$z$是实系数方程$x^2+2bx+c=0$的虚根, 记它在直角坐标平面上的对应点为$P_z(\Re{z},\Im{z})$. 
    若点$(b,c)$在直线$2x+y=0$上, 求证: 点$P_z$在圆$C: (x-1)^2+y^2=1$上. 
    \begin{jieda}
        根据韦达定理有$\left\{\begin{array}{l}
             2Rez = -2b\\
             (Rez)^2 + (Imz)^2 = c
        \end{array}\right.$，因为$2b+c = 0$，整理得到$(Rez-1)^2+(Imz)^2=1$
    \end{jieda}
    
\end{enumerate}


\section{复数集内的方程}
\subsection{复数集内的方程(1)}
\zsd 复数集内的方程 \\
\begin{enumerate}
\item 若方程$x^2+(k+2\mi)x+2+k\mi=0$无实数根, 则实数$k$的取值范围是\blkda[2.5]{$k\neq\pm2\sqrt{2}$}. 


\item 若关于$x$的方程$x^2-kx+2=0$有一个虚根为$\mi$, 则这个方程的另一个根为\blkda{$-2\mi$}. 


\item 已知$m\zc$, 关于$x$的方程$x^2+mx+3+4\mi=0$有实数解, 则复数$m$的模的取值范围是\blkda{$[4,+\infty)$}. 


\item 已知$z\zc$, 关于$x$的方程$x^2-zx+4+3\mi=0$有实数解. 若$\abs{z}=3\sqrt{2}$, 则复数$z=$\blkda[2]{$\pm\frac{3}{\sqrt{5}}(3+\mi)$}. 

\item 方程$3x^2-6(m-1)x+m^2+1=0$的两根均为虚数, 且两根的模的和为2, 则实数$m=$\blkda{$\sqrt{2}$}. 


\item 已知方程$x^5-3x^4+5x^3-5x^2+4x-2=0$有两个根1和$\mi$, 则方程的其他根为\blkda{$1\pm\mi,-\mi$}. 


\item 若$1+\sqrt{2}\mi$是关于$x$的实系数方程$x^2+bx+c=0$的一个复数根, 则\dotfill(\kh{B})
\xx{$b=2,c=3$}{$b=-2,c=3$}{$b=-2,c=-1$}{$b=2,c=-1$}

\item 二次方程$x^2-3\mi{x}-3=0$的根的情况为\dotfill(\kh{B})
\xx{有两个不相等的实根}
{有两个虚根}
{有一对共轭虚根}
{有一个实根和一个虚根}

\item 设$z_1,z_2$是非零复数, 且满足$z_1^2-\sqrt{3}z_1z_2+z_2^2=0$, 则$\abs{z_1},\abs{z_2}$大小的关系是\dotfill(\kh{C})
\xx
{$\abs{z_1}>\abs{z_2}$}
{$\abs{z_1}<\abs{z_2}$}
{$\abs{z_1}=\abs{z_2}$}
{不能确定}


\item 关于一元二次方程$ax^2+bx+c=0(a\neq 0)$的正确结论是\dotfill(\kh{B}) 
\xx
{若$a,b,c\zr$, 则当且仅当$\Delta=b^2-4ac\geq0$时, 方程在复数范围内才有解}
{若$a,b,c\zr$, 则方程在复数范围内总有解}
{若$a,b,c\zc$, 则方程的两根不适合韦达定理}
{若$b=0$时, 则方程的两根互为共轭复数}



\item 解关于$x$的方程$x^2-5x+6+(x-2)\mi=0$. 
\begin{jieda}
    $2$或$3-i$
\end{jieda}


\item 在复数范围内解方程$\abs{z}^2+(z+\overline{z})\mi=\dfrac{3-\mi}{2+\mi}$($\mi$为虚数单位).
\begin{jieda}
    $-\dfrac{1}{2}\pm \dfrac{\sqrt{3}}{2}i$
\end{jieda}

\item 设虚数$z_1,z_2 $满足$z_1^2=z_2$.
\begin{enumerate}[(1)]
\item 若$z_1,z_2$是一个实系数一元二次方程的两个根, 求$z_1,z_2$; 
\item 若$z_1=1+m\mi$($\mi$为虚数单位, $m\zr)$, $\abs{z_1}\leq\sqrt{2}$, 复数$\omega=z_2+3$, 求$\abs{\omega}$的取值范围.
\end{enumerate}
\begin{jieda}
\begin{enumerate}

    \item $z_1=-\dfrac{1}{2}+ \dfrac{\sqrt{3}}{2}i,z_2=-\dfrac{1}{2}-\dfrac{\sqrt{3}}{2}i$；或$z_1=-\dfrac{1}{2}-\dfrac{\sqrt{3}}{2}i,z_2=-\dfrac{1}{2}+\dfrac{\sqrt{3}}{2}i$；
    \item $[\sqrt{13},4)$

\end{enumerate}
\end{jieda}

\item 对任意一个非零复数$z$, 定义集合$M_z=\left\{\omega\mid\omega=z^{2n-1}, n\zn\right\}$. 
\begin{enumerate}[(1)]
\item 设$\alpha$是方程$x+\dfrac{1}{x}=\sqrt{2}$的一个根, 试用列举法表示集合$M_\alpha$; 
\item 设复数$\omega\in M_z$, 求证: $M_\omega\subseteq M_z$.
\end{enumerate}
\begin{jieda}
    (1)$\{\dfrac{\sqrt{2}}{2}(1+i),\dfrac{\sqrt{2}}{2}(-1+i),\dfrac{\sqrt{2}}{2}(1-i),\dfrac{\sqrt{2}}{2}(-1-i)\}$

\end{jieda}

\end{enumerate}


\section{复数的三角形式}
\subsection{复数的三角形式(1)}
\zsd 复数的三角形式; 三角形式下复数的乘除法及几何意义；三角形式下复数的乘方和开方.
A组：

\begin{enumerate}
\item 将下列复数写成代数形式: \\
\NumTabs{2}
\begin{itemize*}[itemjoin={\tab}, label={}]
\item $6\left(\cos\dfrac{3\pi}{2}+\mi\sin\dfrac{3\pi}{2}\right)=$\blkda{$-6\mi$}, 
\item $2\left(\cos\dfrac{5\pi}{3}+\mi\sin\dfrac{5\pi}{3}\right)=$\blkda{$1-\sqrt{3}\mi$}. 
\end{itemize*}

\item  已知$i$是虚数单位，$z=1+\sqrt{3}i$, 则$z$的辅角主值为\blkda{$\dfrac{\pi}{3}$}

\item 已知$i$是虚数单位，$z=-1+\sqrt{3}i$, 则$argz=$\blkda{$\dfrac{2\pi}{3}$}

\item 将下列复数用三角形式表示(用辐角主值): \\
\NumTabs{2} % 
\begin{enumerate*}[(1),itemjoin={\tab}]
\item $4-4\mi=$\blkda[4]{$4\sqrt{2}(\cos\frac{7\pi}{4}+\mi\sin\frac{7\pi}{4})$}; 
\item $-3\sqrt{3}-3\mi=$\blkda[3]{$6(\cos\frac{7\pi}{6}+\mi\sin\frac{7\pi}{6})$}; 
\item $\sin\frac{\pi}{8}+\mi\cos\frac{\pi}{8}=$\blkda[3]{$\cos\frac{3\pi}{8}+\mi\sin\frac{3\pi}{8}$}; 
\item $\overline{\cos\frac{\pi}{7}+\mi\sin\frac{\pi}{7}}=$\blkda[3]{$\cos\frac{13\pi}{7}+\mi\sin\frac{13\pi}{7}$}. 
\end{enumerate*}
\end{enumerate}

B组：
\begin{enumerate}

\item 复数 $z=\sin \theta+i(\cos \theta-1)$ 为纯虚数 $\theta \in[0,2 \pi)$ ，则 $\theta=$\blkda{$\pi$}



\item 将下列复数写成三角形式: \\
\NumTabs{3}
\begin{itemize*}[itemjoin={\tab}, label={}]
\item $-4=$\blkda[3]{$4(\cos\pi+\mi\sin\pi)$}, 
\item $\dfrac{\pi}{2}=$\blkda[3]{$\frac{\pi}{2}(\cos0+\mi\sin0)$}, 
\item $-\pi\mi=$\blkda[3]{$\pi(\cos\frac{3\pi}{2}+\mi\sin\frac{3\pi}{2})$}, 
\item $1-\mi=$\blkda[3.2]{$\sqrt{2}(\cos\frac{7\pi}{4}+\mi\sin\frac{7\pi}{4})$}, 
\item $-1+\sqrt{3}\mi=$\blkda[3]{$2(\cos\frac{2\pi}{3}+\mi\sin\frac{2\pi}{3})$}, 
\item $-\frac{1}{2}-\frac{\sqrt3}{2}\mi=$\blkda[3]{$\cos\frac{4\pi}{3}+\mi\sin\frac{4\pi}{3}$}. 
\end{itemize*}

\item 已知复数 $z_1=-\sqrt{3}+i$ ，若复数 $Z$ 满足 $2 i z=z_1$ ，则复数 $Z$ 的辐角主值为\blkda{$\dfrac{\pi}{3}$}

\item 复数 $\sin 1-i \cos 1$ 的辐角主值是\blkda{$\dfrac{3\pi}{2}+1$}




\NumTabs{2} % 
\item 将下列复数用三角形式表示(用辐角主值): \\
\begin{enumerate*}[(1),itemjoin={\tab}]
\item $\pi\left(\sin1+\mi\cos1\right)=$\blkda[4.5]{$\pi(\cos(\frac{\pi}{2}-1)+\mi\sin(\frac{\pi}{2}-1))$}. 
\item $-2\left(\cos\dfrac{\pi}{4}+\mi\sin\dfrac{\pi}{4}\right)=$\blkda[3]{$2(\cos\frac{5\pi}{4}+\mi\sin\frac{5\pi}{4})$}. 
\item $\cos\dfrac{\pi}{5}-\mi\sin\dfrac{\pi}{5}=$\blkda[3]{$\cos\frac{9\pi}{5}+\mi\sin\frac{9\pi}{5}$}.
\item $-\sqrt{2}\left(\sin\dfrac{\pi}{3}+\mi\cos\dfrac{\pi}{3}\right)=$\blkda[3.5]{$\sqrt{2}(\cos\frac{7\pi}{6}+\mi\sin\frac{7\pi}{6})$}. 
\item $-\dfrac{\pi}{3}\left(\sin\dfrac{\pi}{3}-\mi\cos\dfrac{\pi}{3}\right)=$\blkda[3]{$\frac{\pi}{3}(\cos\frac{5\pi}{6}+\mi\sin\frac{5\pi}{6})$}.
\item $-2\left(\cos\dfrac{\pi}{3}-\mi\sin\dfrac{\pi}{3}\right)=$\blkda[3]{$2(\cos\frac{2\pi}{3}+\mi\sin\frac{2\pi}{3})$}. 
\item $\sqrt{3}\left(\sin\dfrac{\pi}{6}-\mi\cos\dfrac{\pi}{6}\right)=$\blkda[3.2]{$\sqrt{3}(\cos\frac{5\pi}{3}+\mi\sin\frac{5\pi}{3})$}. 
\end{enumerate*}

\item 若 $z=\cos \theta+i \sin \theta$（ $i$ 为虚数单位），则 $z^2=-1$ 的 $\theta$ 值可能是\blkda{$\dfrac{\pi}{2}$}

\item 设复数 $z$ 满足条件 $\arg z \in\left(\frac{3}{4} \pi, \pi\right)$ ，则 $\frac{2021}{\overline{z^2}}$ 对应复平面上的点位于第 \blkda{四}象限．

\item 若 $\left(\cos \frac{\pi}{8}+i \sin \frac{\pi}{8}\right)^n$ 是纯虚数（其中 $i$ 是虚数单位），则正整数 $n$ 的最小值为 \blkda{4}．

\item $i$ 的所有的三次方根为 \blkda{$-1,\pm\dfrac{\sqrt{3}}{2}+\dfrac{1}{2}i$} ．

\item 已知复数 $z$ 满足 $|z|=2, \arg (z+2)=\frac{\pi}{3}$ ，则 $z$ 为 \blkda{$-1+\sqrt{3}i$} ．

\item 方程 $z^7=-1+\sqrt{3} i$ 的 7 个根在复平面上对应了 7 个点，第\blkda{三}象限只包含了其中的一个点．



\item 已知复数 $z_1=2 \sin \theta-\sqrt{3} i, z_2=1+(2 \cos \theta) i, i$ 为虚数单位，若 $\theta \in\left[\frac{\pi}{2}, \frac{5 \pi}{6}\right]$ ，复数 $z_1, z_2$ 对应的向量分别为 $\vec{a}, \vec{b}$ ，存在 $\theta$ 使得等式 $(\lambda \vec{a}-\vec{b}) \cdot(\vec{a}-\lambda \vec{b})=0$ 成立，则实数 $\lambda$ 的取值范围为 \blkda[4]{$[2-\sqrt{3},2+\sqrt{3}]$}．

\item 在复平面上，设点 $A 、 B$ 对应的复数分别为 $-2 i 、 \cos \left(t-\frac{\pi}{3}\right)+i \sin \left(t-\frac{\pi}{3}\right)$ ，当 $t$ 由 $\frac{\pi}{6}$ 连续变到 $\frac{\pi}{2}$ 时，向量 $\overrightarrow{A B}$ 所扫过的图形区域的面积是 \blkda{$\dfrac{\pi}{6}$}．



\item 计算, 并将结果写成代数形式: 
\begin{enumerate}[(1)]
\item $8\left(\cos\dfrac{\pi}{6}+\mi\sin\dfrac{\pi}{6}\right)\cdot2\left(\cos\dfrac{\pi}{3}+\mi\sin\dfrac{\pi}{3}\right)=$\blkda{$16\mi$}.  

\item $\dfrac{6\left(\cos\frac{3\pi}{5}+\mi\sin\frac{3\pi}{5}\right)}{\sqrt{2}\left(\cos\frac{\pi}{10}+\mi\sin\frac{\pi}{10}\right)}=$\blkda{$3\sqrt{2}\mi$}. 

\item $3\left(\cos\ang{18}+\mi\sin\ang{18}\right)\times 2\left(\cos\ang{54}+\mi\sin\ang{54}\right)\times 5\left(\cos\ang{108}+\mi\sin\ang{108}\right)=$\blkda{$-30$}. 

\item $\sqrt{3}\left(\cos\ang{150}+\mi\sin\ang{150}\right)\div\left[\sqrt{2}\left(\cos\ang{300}+\mi\sin\ang{300}\right)\right]=$\blkda[3]{$-\frac{3\sqrt{2}}{4}-\frac{\sqrt{6}}{4}\mi$}. 

\item $\left(\cos\dfrac{\pi}{4}-\mi\sin\dfrac{\pi}{4}\right)^3=$\blkda[3]{$-\frac{\sqrt{2}}{2}-\frac{\sqrt{2}}{2}\mi$}. 

\item $1\div\left(\cos\dfrac{\pi}{4}+\mi\sin\dfrac{\pi}{4}\right)=$\blkda[3]{$\frac{\sqrt{2}}{2}-\frac{\sqrt{2}}{2}\mi$}. 
\end{enumerate}



\item 在复数域中解方程：
（1）$x^4=1$
（2）$\omega^3=-i$
（3）$x^3=1-i$
（4）$x^2=1-\sqrt{3} i$


\item 已知 $f(z)=\bar{z}-1$ ，且 $f\left(z_1-z_2\right)=4+4 i$ ，若 $z_1=2-2 i$ ．
\begin{enumerate}
    \item 求复数 $z_1=2-2 i$ 的三角形式，并且复数 $Z_1$ 的辐角主值 $\arg z_1$ ；
    \item 求 $\left|\frac{z_1-z_2}{z_1+z_2}\right|$ ．
\end{enumerate}

\begin{jieda}
    (1)$z_1=2\sqrt{2}(cos\dfrac{7\pi}{4})+isin\dfrac{7\pi}{4}$, argz_1=\dfrac{7\pi}{4};(2)\sqrt{41}
\end{jieda}



\item 已知复数 $z$ 的模为 1 ，辐角主值为 $\theta, \theta \in(0,2 \pi)$ ，试用 $\theta$ 表示复数 $\omega=z^2-z$ 的模及辐角主值．

\begin{jieda}
    $|\omega|=2sin\dfrac{\theta}{2}$;
    $arg(z-1)z = \left \{ \begin{array}{ll}
                     \frac{\pi+3\theta}{2}, &  \theta\in [0,\pi) \\
                     \frac{3\theta}{2}-\frac{3\pi}{2}, & \theta\in [\pi, 2\pi) \\
                \end{array} \right.$
\end{jieda}

\item 设复数 $z_1=\sqrt{3}+i$ ，复数 $z_2$ 满足 $\left|z_2\right|=2$ ，已知 $z_1 \cdot z_2^2$ 的对应点在虚轴的负半轴上，且 $\arg z_2 \in(0, \pi)$ ，求 $z_2$ 的代数形式．
\begin{jieda}
    $-1+\sqrt{3}i$
\end{jieda}

\item 已知复数 $z_1$ 的辐角主值是 $\frac{5 \pi}{12}$ ，且有 $\left|z_1-z_0\right|=\sqrt{2}, z_0=(1+i) z_1$ ．
（1）求复数 $z_0$ 的代数表示；
（2）求满足 $\left|z-z_0\right|=\sqrt{2}$ 的复数 $z$ 的辐角主值的最大值．
\begin{jieda}
    (1) $-1+\sqrt{3}i$;(2)$\dfrac{11\pi}{12}$
\end{jieda}




\end{enumerate}